\subsection{Discrete Model}
Define the force of infection in region $i$ at time $t$:
\begin{equation*}
\lambda(t) = \beta(t) \mathbf{C}I(t)
\end{equation*}

Over one time step, the probability that a susceptible individual in
region $i$ does \emph{not} become infected is:
\begin{equation*}
\exp(-\lambda_i(t))
\end{equation*}

and the probability of infection is:
\begin{equation*}
p_i(t) = 1 - \exp(-\lambda_i(t))
\end{equation*}

The expected incidence:
\begin{equation*}
\mu_i(t) = S_i(t) \cdot p_i(t) = S_i(t) \left[1 - \exp(-\lambda_i(t))\right]
\end{equation*}

The susceptible population evolves as:
\begin{equation*}
S_i(t+1) = S_i(t) - \mu_i(t) = S_i(t) [1 - p_i(t)] = S_i(t) \exp(-\lambda_i(t))
\end{equation*}


The probability an infected individual remains infected over one time
step is:
\begin{equation*}
\exp(-\gamma)
\end{equation*}

The probability of recovery is:
\begin{equation*}
q = 1 - \exp(-\gamma)
\end{equation*}

Recalling the incidence is $\mu(t) = S(t) [1 - \exp(-\lambda(t))]$:
\begin{align}\begin{split}\label{eq:exp_sir_compact}
S(t+1) &= S(t) \exp(-\lambda(t)) \\
I(t+1) &= I(t) \exp(-\gamma) + \mu(t) \\
R(t+1) &= R(t) + I(t) [1 - \exp(-\gamma)]
\end{split}\end{align}


\subsection{Parameters}

\textbf{Derivation of $\gamma$:} We aim to match a continuous model
with infectious period $2.2$ days using a discrete weekly time step
model. We match the \emph{decay rate} over one time step. In the continuous
model with recovery rate $\gamma_{\text{cont}} = 1/2.2$ per day, the
infected fraction after 7 days is $\exp(-7/2.2)$. In our discrete
model, after 1 week it is $\exp(-\gamma)$. Matching these:
$$
\gamma = \frac{7}{2.2} \approx 3.18
$$
Mean infectious duration is $1/q \approx 1.04$ weeks.

\textbf{Derivation of $R_0$ and $\beta$:} Consider population size $N$ (large) with one infected individual initially. Ignoring seasonality and contact structure:
$$\lambda_i(t) = \frac{\beta}{N}$$

The infection probability per timestep:
$$p_i(t) = 1 - \exp(-\lambda_i(t)) = 1 - \exp(-\beta/N)$$

For large $N$, using $\exp(-x) \approx 1-x$ for small $x$:
$$p_i(t) \approx 1 - (1 - \beta/N) = \frac{\beta}{N}$$

Incidence per timestep:
$$\mu_i(t) = S_i(t) \cdot p_i(t) = (N-1) \cdot \frac{\beta}{N} \to \beta$$
as $N \to \infty$.

The infected individual recovers with probability $q = 1-\exp(-\gamma)$ per timestep. Mean infectious duration is $1/q$ timesteps. Total secondary infections:
$$R_0 = \beta \times \frac{1}{q} = \frac{\beta}{1-\exp(-\gamma)}$$

With seasonal forcing, transmission varies as:
$$\beta(t) = \beta_0\left(1 + A\cos\left(\frac{2\pi t}{P}\right)\right)$$
where $A$ is the amplitude and $P$ is the period.

The measured $R_0 = 3.65$ corresponds to peak transmission when $\cos(2\pi t/P) = 1$:
$$R_0 = \frac{\beta_0(1 + A)}{1-\exp(-\gamma)}$$

Therefore:
$$\beta_0 = \frac{R_0 \left(1 - \exp(-\gamma)\right)}{1 + A}$$

For example, with $R_0 = 3.65$, $\gamma = 7/2.2$, and amplitude $A = 0.5$:
$$\beta_0 = \frac{3.65 \times 0.958}{1.5} \approx 2.33$$

%% However, attempting to match mean durations directly is
%% \emph{mathematically impossible}.

%% To see why: if we try to match the mean duration $1/(1-\exp(-\gamma)) = 2.2/7 \approx 0.314$ weeks, this would require:
%% $$1 - \exp(-\gamma) = 7/2.2 \approx 3.18$$
%% $$\exp(-\gamma) = -2.18$$

%% This is impossible since $\exp(-\gamma) > 0$ always. The fundamental problem is that with weekly time steps, we cannot accurately represent a $2.2$-day infectious period ($< 1$ week).



\subsection{Sensitivities etc.}
Taking partials:
\begin{align*}
  \frac{\partial \lambda(t)}{\partial \phi_k} &= \beta(t) \left
  (\frac{\partial \mathbf{C}}{\partial \phi_k}I(t) + \mathbf{C}
  \frac{\partial I(t)}{\partial \phi_k} \right ) \\
  %
  %
  %
  \frac{\partial \mu(t)}{\partial \phi_k} &= \frac{\partial
    S(t)}{\partial \phi_k} \left[1 - \exp(-\lambda(t))\right] + S(t)
  \exp(-\lambda(t)) \frac{\partial \lambda(t)}{\partial \phi_k}
\end{align*}
%
%
where
%
%
\begin{equation*}
\frac{\partial \mathbf{C}}{\partial \phi_k} = \begin{cases}
\begin{bmatrix} -1 & 1 \\ 1 & -1 \end{bmatrix} & \phi_k = \theta \\
0 &  \text{o.w.}
\end{cases}
\end{equation*}
%
%
Taking partials of the dynamics:
%
%
\begin{align}\begin{split}
    \frac{\partial S(t+1)}{\partial \phi_k}
    %
    %
    %
    &= \exp(-\lambda(t)) \frac{\partial S(t)}{\partial \phi_k} -
    S(t) \exp(-\lambda(t)) \frac{\partial \lambda(t)}{\partial
      \phi_k} \\
    %
    %
    %
    %% &= \exp(-\lambda_i(t)) \frac{\partial S_i(t)}{\partial \phi_k} -
    %% S_i(t) \exp(-\lambda_i(t)) \frac{\partial \lambda_i(t)}{\partial
    %%   \phi_k}\\
    %
    %
    %
    \frac{\partial I(t+1)}{\partial \phi_k} &= \exp(-\gamma)
    \frac{\partial I(t)}{\partial \phi_k} + \frac{\partial
      \mu(t)}{\partial \phi_k}
\end{split}\end{align}


We assume a constant reporting rate $\rho$. Observations are binomial
\begin{align}\begin{split}
Y_i(t) &\sim Bin(\mu_i(t), \rho) \\
\mathbb{E}[Y_i(t)] &= \rho \mu_i(t) \\
\mathbb{V}ar[Y_i(t)] &= \rho(1-\rho) \mu_i(t) \\
\end{split}\end{align}

We can approximate this using the central limit theorem as a gaussian
\begin{align*}
  Y_i(t) &= \rho \mu_i(t) + \epsilon_i(t) \\
  \epsilon_i(t) &\sim \mathcal{N}(0, \rho(1-\rho)\mu_i(t))
\end{align*}

The log-likelihood for independent Gaussian observations is:
\begin{equation*}
\ell(\boldsymbol{\phi}) = -\frac{1}{2}\sum_{t=0}^{T-1}\sum_{i=1}^{2} \left[\log(2\pi \rho(1-\rho)\mu_i(t)) + \frac{(Y_i(t) - \rho\mu_i(t))^2}{\rho(1-\rho)\mu_i(t)}\right]
\end{equation*}

Taking the partial derivative with respect to parameter $\phi_k$:
\begin{align*}
\frac{\partial \ell}{\partial \phi_k} &= -\frac{1}{2}\sum_{t,i} \left[\frac{1}{\mu_i(t)} - \frac{2\rho(Y_i(t) - \rho\mu_i(t))}{\rho(1-\rho)\mu_i(t)} - \frac{(Y_i(t) - \rho\mu_i(t))^2}{\rho(1-\rho)\mu_i(t)^2}\right] \frac{\partial \mu_i(t)}{\partial \phi_k}
\end{align*}

We can write this as:
\begin{equation*}
\frac{\partial \ell}{\partial \phi_k} = \sum_{t,i} A_i(t) \frac{\partial \mu_i(t)}{\partial \phi_k}
\end{equation*}
where
\begin{equation*}
A_i(t) = -\frac{1}{2\mu_i(t)} + \frac{Y_i(t) - \rho\mu_i(t)}{(1-\rho)\mu_i(t)} + \frac{(Y_i(t) - \rho\mu_i(t))^2}{2\rho(1-\rho)\mu_i(t)^2}
\end{equation*}

The Fisher Information Matrix is:
\begin{equation*}
[\mathbf{J}]_{kl} = \mathbb{E}\left[\frac{\partial \ell}{\partial \phi_k}\frac{\partial \ell}{\partial \phi_l}\right] = \sum_{t,i}\sum_{s,j} \mathbb{E}[A_i(t) A_j(s)] \frac{\partial \mu_i(t)}{\partial \phi_k}\frac{\partial \mu_j(s)}{\partial \phi_l}
\end{equation*}

Since observations are independent across time and regions, $\mathbb{E}[A_i(t) A_j(s)] = 0$ when $(i,t) \neq (j,s)$. For $(i,t) = (j,s)$, we need $\mathbb{E}[A_i(t)^2]$.

Recalling that $Y_i(t) - \rho\mu_i(t) = \epsilon_i(t)$, we can write:
\begin{equation*}
A_i(t) = -\frac{1}{2\mu_i(t)} + \frac{\epsilon_i(t)}{(1-\rho)\mu_i(t)} + \frac{\epsilon_i(t)^2}{2\rho(1-\rho)\mu_i(t)^2}
\end{equation*}

Expanding $A_i(t)^2$:
\begin{align*}
A_i(t)^2 &= \frac{1}{4\mu_i(t)^2} - \frac{\epsilon_i(t)}{(1-\rho)\mu_i(t)^2} - \frac{\epsilon_i(t)^2}{2\rho(1-\rho)\mu_i(t)^3} \\
&\quad + \frac{\epsilon_i(t)^2}{(1-\rho)^2\mu_i(t)^2} + \frac{\epsilon_i(t)^3}{\rho(1-\rho)^2\mu_i(t)^3} + \frac{\epsilon_i(t)^4}{4\rho^2(1-\rho)^2\mu_i(t)^4}
\end{align*}

Using the moments of $\epsilon_i(t) \sim \mathcal{N}(0, \rho(1-\rho)\mu_i(t))$:
\begin{itemize}
\item $\mathbb{E}[\epsilon_i(t)] = 0$
\item $\mathbb{E}[\epsilon_i(t)^2] = \rho(1-\rho)\mu_i(t)$
\item $\mathbb{E}[\epsilon_i(t)^3] = 0$ (third moment of centered Gaussian)
\item $\mathbb{E}[\epsilon_i(t)^4] = 3[\rho(1-\rho)\mu_i(t)]^2$ (fourth moment of Gaussian)
\end{itemize}

Taking expectations:
\begin{align*}
\mathbb{E}[A_i(t)^2] &= \frac{1}{4\mu_i(t)^2} - \frac{\rho(1-\rho)\mu_i(t)}{2\rho(1-\rho)\mu_i(t)^3} + \frac{\rho(1-\rho)\mu_i(t)}{(1-\rho)^2\mu_i(t)^2} \\
&\quad + \frac{3[\rho(1-\rho)\mu_i(t)]^2}{4\rho^2(1-\rho)^2\mu_i(t)^4} \\
&= \frac{1}{4\mu_i(t)^2} - \frac{1}{2\mu_i(t)^2} + \frac{\rho}{(1-\rho)\mu_i(t)^2} + \frac{3}{4\mu_i(t)^2} \\
&= \frac{1}{2\mu_i(t)^2} + \frac{\rho}{(1-\rho)\mu_i(t)^2}\\
%
%
%
&= \frac{1+\rho}{2(1-\rho)\mu_i(t)^2}
\end{align*}

Therefore:
\begin{equation*}
[\mathbf{J}]_{kl} = \frac{1+\rho}{2(1-\rho)}\sum_{t=0}^{T-1}\sum_{i=1}^{2} \frac{1}{\mu_i(t)^2} \frac{\partial \mu_i(t)}{\partial \phi_k}\frac{\partial \mu_i(t)}{\partial \phi_l}
\end{equation*}

In matrix form, define the Jacobian matrix:
\begin{equation*}
\mathbf{G} = \frac{\partial \boldsymbol{\mu}}{\partial \boldsymbol{\phi}} \in \mathbb{R}^{2T \times 5}
\end{equation*}
where $G_{(t,i),k} = \frac{\partial \mu_i(t)}{\partial \phi_k}$.

Define the diagonal weight matrix:
\begin{equation*}
\mathbf{W} = (\frac{1+\rho}{2(1-\rho)} \text{diag}\left( \mu_1(0)^{-2}, \ldots, \mu_2(T-1)^{-2}\right) \in \mathbb{R}^{2T \times 2T}
\end{equation*}

Then the Fisher Information Matrix is:
\begin{equation*}
\mathbf{J} = \mathbf{G}^T \mathbf{W} \mathbf{G}
\end{equation*}

\subsubsection{Cramér-Rao Lower Bound}

The Cramér-Rao lower bound for the connectivity parameter $\theta$ is:
\begin{equation*}
\text{Var}(\hat{\theta}) \geq [\mathbf{J}^{-1}]_{11}
\end{equation*}

This can be computed efficiently using Cholesky factorization. Let $\mathbf{J} = \mathbf{L}\mathbf{L}^T$, then:
\begin{equation*}
[\mathbf{J}^{-1}]_{11} = \|\mathbf{L}^{-T}\mathbf{e}_1\|^2
\end{equation*}
where $\mathbf{e}_1 = [1, 0, 0, 0, 0]^T$ is the first standard basis vector.

The CRLB for the standard deviation of $\hat{\theta}$ is:
\begin{equation*}
\text{SD}(\hat{\theta}) \geq \sqrt{[\mathbf{J}^{-1}]_{11}} = \|\mathbf{L}^{-T}\mathbf{e}_1\|
\end{equation*}
